\section{Resultados e discussão}
\label{sec:resultados}

Apresente os resultados na mesma ordem das etapas do método. Antes de cada
tabela ou figura, explique o que o leitor deve observar. Depois, interprete o
resultado, compare-o com a literatura e discuta implicações e limitações.

% Exemplo de figura. Coloque o arquivo em figuras/ e remova o símbolo %.
% \begin{figure}[htbp]
%   \centering
%   \includegraphics[width=0.78\textwidth]{nome-da-figura.png}
%   \caption{Título informativo da figura.}
%   \label{fig:resultado-principal}
% \end{figure}

% Exemplo de tabela editável.
\begin{table}[htbp]
  \centering
  \caption{Exemplo de comparação entre cenários.}
  \label{tab:cenarios}
  \begin{tabular}{lcc}
    \toprule
    Cenário & Indicador 1 & Indicador 2 \\
    \midrule
    Referência & -- & -- \\
    Alternativo & -- & -- \\
    \bottomrule
  \end{tabular}
\end{table}

Substitua a tabela de exemplo e evite repetir no texto todos os valores que já
estão visíveis. Destaque padrões, diferenças e respostas à pergunta de pesquisa.
